\subsection{An approximation bound and a discarded Dedekind construction}

\begin{lemma}[Proved: approximation by finitely generated subalgebras]
\label{C:A2-approximation}
Let $k$ be a field, let $\ell\ne\operatorname{char}k$, and suppose
$c=\operatorname{cd}_{\ell}k<\infty$.  Let $R$ be a $k$-algebra domain
whose fraction field $K$ has finite transcendence degree $t$ over $k$.
We do not assume that $R$ is noetherian or finitely generated.  Then
\[
\operatorname{cd}_{\ell}(\operatorname{Spec}R)\le c+t.
\]
In particular, if $k$ is algebraically closed, $K$ is algebraic over
$k(t)$, and $R$ is a Dedekind subring of $K$ with fraction field $K$,
then $R$ contains $k$, its $\ell$-cohomological dimension is at most
one, and $\operatorname{Pic}(R)/\ell=0$.
\end{lemma}

\begin{proof}
Write $R=\varinjlim_i R_i$, where $R_i$ runs through its finitely
generated $k$-subalgebras.  Every $R_i$ is a domain and
$\dim R_i=\operatorname{trdeg}_k\operatorname{Frac}(R_i)\le t$.
Artin's affine theorem over a separable closure of $k$, followed by
the Hochschild--Serre spectral sequence, gives
$\operatorname{cd}_{\ell}(\operatorname{Spec}R_i)\le c+t$; the
precise affine theorem is SGA~4, Expos\'e~XIV, Corollary~3.2
\cite{C_SGA4}.

Here are the two limit steps, including nonconstructible coefficients.
Every constructible $\mathbf Z/\ell^n$-sheaf on
$\operatorname{Spec}R$ descends to some $\operatorname{Spec}R_i$,
and its cohomology is the filtered colimit of the cohomology of its
pullbacks to later stages.  Consequently it has no cohomology in
degrees above $c+t$.  Every $\ell$-primary torsion sheaf is a filtered
colimit of constructible $\ell$-primary torsion sheaves, and cohomology
on the affine scheme $\operatorname{Spec}R$ commutes with that
filtered colimit.  We use, respectively, the Stacks Project,
Lemma~59.73.10(3), Theorem~59.51.3, Lemma~59.73.2(2)--(3), and
Lemma~59.51.4, all opened in this run \cite{C_StacksLimits}.
This proves the assertion for arbitrary torsion coefficients.

For the final assertion, consider a maximal ideal $\mathfrak m$ of
the Dedekind domain $R$ and its normalized discrete valuation
$v_{\mathfrak m}\colon K^\times\to\mathbf Z$.  The group $k^\times$
is divisible because $k$ is algebraically closed.  Its image in
$\mathbf Z$ is therefore zero.  Thus every nonzero element of $k$
is a unit in every $R_{\mathfrak m}$.  The intersection identity
$R=\bigcap_{\mathfrak m}R_{\mathfrak m}$ inside $K$ gives $k\subset R$.
If $R$ is itself a field, this inclusion is immediate.  We may now
apply the first assertion with $c=0$ and $t=1$.  Finally the Kummer
sequence injects $\operatorname{Pic}(R)/\ell$ into
$H^2(R,\mu_\ell)=0$.
\end{proof}

This rules out a persistent Picard-group obstruction obtained by
prescribing a Dedekind class group while keeping the fraction field
algebraic over a one-variable function field over an algebraically
closed field.  It does not rule out algebraic extensions of number
fields: there the cohomological dimension of the fraction field can
drop below the bound furnished by the original ground field.

\begin{lemma}[Proved: finite character in an algebraic tower]
\label{C:A2-tower-Dedekind}
Let $R_0$ be a Dedekind domain and let $K_\infty/K_0$, with
$K_0=\operatorname{Frac}(R_0)$, be an algebraic extension expressed as
a directed union of finite extensions $K_i$.  Suppose that their
integral closures $R_i$ of $R_0$ are finite over $R_0$.  Suppose also
that each nonzero prime of every $R_i$ has only finitely many primes
above it in $R_\infty=\bigcup_iR_i$, and that along each such prime
the ramification indices in the tower are bounded.  Then
$R_\infty$ is a Dedekind domain.
\end{lemma}

\begin{proof}
Integrality shows that every nonzero prime of $R_\infty$ is maximal.
At such a prime $P$, the localization $(R_\infty)_P$ is the directed
union of the compatible valuation rings $(R_i)_{P\cap R_i}$.  It is
a valuation ring.  Bounded ramification indices make its value group
a discrete cyclic subgroup of $\mathbf Q$, so this localization is
a discrete valuation ring.  A nonzero element $a\in R_\infty$
belongs to some $R_i$ and has only finitely many prime divisors in
that $R_i$; the finite-splitting hypothesis then gives only finitely
many maximal ideals of $R_\infty$ containing $a$.

We check noetherianity directly.  Given a nonzero ideal $I$, choose
$0\ne a\in I$, and list the finitely many maximal ideals
$P_1,\dots,P_s$ containing $a$.  For each $j$ choose $b_j\in I$
whose valuation at $P_j$ is the minimum of the nonnegative integer
valuations attained on $I\setminus\{0\}$.  The ideal
$J=(a,b_1,\dots,b_s)$ has $J_{P_j}=I_{P_j}$ for every $j$; at every
other maximal ideal $a$ is a unit, so again $J_P=I_P$.  Local equality
of these ideals at all maximal ideals implies $I=J$.  Thus every
ideal is finitely generated.  The domain is noetherian, normal and
one-dimensional, as required.
\end{proof}
